% ============================================================
% Discussion --- interpretation, implications, limitations (target: C&S)
% ============================================================
\section{Discussion}
\label{sec:discussion}

The results converge on a single thesis: residual spear phishing is detectable
exactly to the degree to which it \emph{violates} the victim's observable domain knowledge ---
and that degree can be \emph{measured}. Below we interpret the mechanism, outline the
practical implications, and honestly delineate the limits.

\subsection{Interpretation: two invariants, two failure modes}
The two invariants are not competing detectors but \emph{complementary} responses to two
disjoint modes of targeted attack. The cross-layer inconsistency invariant
(\ref{sec:res-inv1}) addresses the \emph{new sender}: role or authority impersonation is
``known'' in one layer (e.g.\ name, domain), yet inconsistent in others (no OSINT trace, no
relationship history) --- which breaks the pattern of a genuine relationship that is consistent across many layers at once.
The cascade dynamics invariant (\ref{sec:res-inv2}) addresses the \emph{compromised account}, which is
consistent across all static layers (because it \emph{is} a genuine contact), yet betrays itself
\emph{temporally}: it sends out of rhythm, creating a recipient-side cascade. The per-threat
ablation makes this complementarity explicit: for new senders the signal is carried by
OSINT, while time is useless; for compromises it is the reverse. This observation explains
why single approaches (a content filter alone, a contact graph alone, a volumetric detector alone) fail
on the residual $2\%$: each is blind to one of the two modes.

Mechanistically, the most important result concerns the \emph{stealthy attacker}. A volumetric detector
(in the spirit of COMPA \citep{stringhini2015thataint}) assumes that a compromise manifests as a ``surge'';
when the attacker deliberately spreads out the sending, the volumetric signal vanishes (AUC below random).
Cascade dynamics reads the \emph{recipient side} --- the wave-like signature among the neighbors ---
whose stealth the sender cannot hide. This is consistent with the observations of Ho et al.\ \citep{ho2019detecting},
that lateral phishing requires a signal beyond the content of a single message; we differ, however,
in that we locate the signal in the \emph{dynamics of the domain knowledge graph}, not in headers or relays.

\subsection{Practical implications}
The strongest implication stems from the susceptibility measurement (\ref{sec:susceptibility}) and is \emph{counter-intuitive}.
Typical defense programs prioritize ``VIPs'' (executives) as the most exposed. Our
measurement shows that with respect to the \emph{lateral} channel the opposite holds: executives, densely embedded
in the domain knowledge graph, are \emph{best} protected by the structural signal alone, while the largest
unaddressed margin lies on the \emph{periphery} (mid/junior). The recommendation is therefore two-part:
(a)~awareness programs and hard controls (e.g.\ out-of-band verification for transfers) should
cover the \emph{periphery}, not just the top; (b)~detection resources are worth directing where the domain knowledge
graph is sparse --- because that is where the residual margin is largest. With respect to the
\emph{cold impersonation} channel the measurement delivers a hard message: no relational layer suffices
($\mathrm{Susc}\!\approx\!0.96$ regardless of rank), so here the graph \emph{must} be combined
with content and process controls --- as the fusion experiment confirms. This makes the domain knowledge graph
not a standalone product but a \emph{prioritization layer}: it indicates who, and through which relationship,
is most weakly protected.

The result also has an implication for the human factor. Because human susceptibility rises under
time pressure and authority cues while awareness training moderates it only partially
(Section~\ref{sec:related}), a structural susceptibility measure allows limited training interventions
to be \emph{targeted} at the roles that the technology protects most weakly.

\subsection{Limitations}
We treat the limitations seriously, because they define the scope of the claims. \emph{First,
maliciousness is synthetic.} No public corpus links real OSINT facts with labeled compromises;
we generate attack labels on real structure (Enron, SNAP) or on a synthetic population of twins.
We therefore validate the \emph{mechanism} (out-of-rhythm mail separates compromises; layer
inconsistency separates impersonation), not the detection of real intrusions. The OSINT multiplex
advantage (Invariant~I) remains \textbf{exclusively synthetic} --- we read it
as a proof of soundness, not a measurement on real traffic. We also note a degree of
\emph{circularity}: the synthetic population is generated graph-first, so detecting a violation of that
structure is in part built in. Consistent with this, the layer analysis (\ref{sec:res-inv1}) shows the
\emph{temporal} layer carries almost all of Invariant~I's signal, the structural and OSINT layers
contributing little unless the topology is artificially decorrelated; the headline of Invariant~I thus
reduces, in practice, to ``rhythm consistency detects off-rhythm mail'', whereas the cascade result
(Invariant~II) establishes the receiving-side signal far more robustly on real graphs. \emph{Second, the susceptibility measurement
inherits these assumptions}: the $\mathrm{Susc}$ figures are relative to our event model
and population structure (19~companies~$\times\sim$8~people, one domain per company); the direction of the conclusions
(relationship~$>$~rank; periphery most exposed) is qualitatively robust, but the exact values
are not directly transferable to every organization. The direction is moreover partly
\emph{definitional}: executives are layer-rich \emph{by construction}, so their stronger structural
protection follows in part from how we tied layer richness to seniority --- a coupling that should be
stress-tested under alternative assignments before the ``periphery first'' recommendation is generalized. \emph{Third, the scale is limited}
($N{=}150$ twins; real graphs up to~$\sim$2k nodes), so the comparison
with heavy GNN architectures (GATNE/HGT) is treated as a sanity check, not a full benchmark.
\emph{Fourth, the operational threshold}: even the best operating point (sensitivity@FPR$1\%$) remains
orders of magnitude away from the deployment threshold on real mail (Ho et al.\ \citep{ho2019detecting}:
$<\!4$~FP/M); we report it as a \emph{requirement we are beginning to meet}, not an absolute record.

\subsection{Future work directions}
\begin{itemize}[noitemsep,topsep=2pt]
  \item \textbf{Real BEC labels.} The most important next step is validation on a corpus
  with labeled, real compromises --- to turn the proof of mechanism into a measurement of detectability.
  \item \textbf{A learned, inductive domain knowledge graph.} Replacing the hand-crafted consistency features
  with a learned multiplex GNN with explicit cross-layer consistency regularization, validated
  inductively on unseen organizations.
  \item \textbf{Adaptive attack and drift.} Studying an attacker that \emph{learns}
  the domain knowledge graph (LLM~$+$~OSINT \citep{bethany2024evaluating}) and tunes rhythm/relationships, as well as population drift
  over time (TESSERACT \citep{pendlebury2019tesseract}).
  \item \textbf{Susceptibility as a steering signal.} Using the measured $\mathrm{Susc}$ for
  dynamic prioritization of controls and personalization of training --- with an evaluation of whether it lowers the real
  click-through rate.
\end{itemize}
